\documentclass[11pt,a4paper]{article}

\usepackage[utf8]{inputenc}
\usepackage[T1]{fontenc}
\usepackage[french]{babel}
\usepackage[a4paper,margin=2cm]{geometry}
\usepackage{amsmath,amssymb,amsthm,mathtools}
\usepackage{lmodern}
\usepackage{enumitem}
\usepackage{fancyhdr}
\usepackage{lastpage}
\usepackage{xcolor}
\usepackage{array}
\usepackage{tabularx}

\setlength{\headheight}{14pt}
\pagestyle{fancy}
\fancyhf{}
\fancyhead[L]{Terminale générale -- Spécialité Mathématiques + Maths expertes}
\fancyhead[C]{Corrigé détaillé complet}
\fancyhead[R]{Page \thepage/\pageref{LastPage}}
\fancyfoot[C]{\textit{Corrigé détaillé avec rappels, méthodes et commentaires}}

\newcommand{\R}{\mathbb{R}}
\newcommand{\C}{\mathbb{C}}
\newcommand{\Z}{\mathbb{Z}}
\newcommand{\N}{\mathbb{N}}
\newcommand{\e}{\mathrm{e}}
\newcommand{\ii}{\mathrm{i}}

\newtheorem*{rappel}{Rappel}
\newtheorem*{methode}{Méthode}
\newtheorem*{remarque}{Remarque}
\newtheorem*{reponseattendue}{Réponse attendue}

\begin{document}

\begin{center}
    {\Large \textbf{CORRIGÉ DÉTAILLÉ -- MATHÉMATIQUES}}\\[0.3em]
    {\large \textbf{Terminale générale -- Spécialité Mathématiques + option Maths expertes}}\\[0.4em]
    {\large \textbf{Sujet d’approfondissement -- version corrigée}}\\[0.5em]
    Avec rappels de méthodes, commentaires et réponses attendues
\end{center}

\vspace{0.5em}
\hrule
\vspace{0.6em}

\section*{Exercice 1 -- Fonctions, inégalités, raisonnement global}

On considère la fonction
\[
g(x)=x-\ln x
\qquad \text{définie sur } ]0,+\infty[.
\]

\subsection*{Partie A -- Étude de \(g\)}

\begin{enumerate}[label=\textbf{1.\arabic*}, leftmargin=1.6em]

\item \textbf{Limites de \(g\)}

Quand \(x\to 0^+\), on sait que
\[
\ln x \to -\infty.
\]
Donc
\[
g(x)=x-\ln x \to 0-(-\infty)=+\infty.
\]

Quand \(x\to +\infty\), on sait que \(x\) domine \(\ln x\), donc
\[
x-\ln x \to +\infty.
\]

\[
\boxed{\lim_{x\to 0^+} g(x)=+\infty
\qquad\text{et}\qquad
\lim_{x\to +\infty} g(x)=+\infty.}
\]

\item \textbf{Variations de \(g\)}

\[
g'(x)=1-\frac1x=\frac{x-1}{x}.
\]

Comme \(x>0\), le signe de \(g'(x)\) est celui de \(x-1\). Donc :
\begin{itemize}
    \item \(g'(x)<0\) sur \(]0,1[\),
    \item \(g'(1)=0\),
    \item \(g'(x)>0\) sur \(]1,+\infty[\).
\end{itemize}

Ainsi, \(g\) est décroissante sur \(]0,1]\), puis croissante sur \([1,+\infty[\).

\item \textbf{Montrer que \(x-\ln x\ge 1\)}

D’après la question précédente, \(g\) admet un minimum en \(x=1\). Or
\[
g(1)=1-\ln 1=1.
\]

Donc, pour tout \(x>0\),
\[
g(x)\ge g(1)=1,
\]
soit
\[
\boxed{x-\ln x\ge 1.}
\]

Le cas d’égalité a lieu si et seulement si
\[
\boxed{x=1.}
\]

\end{enumerate}

\subsection*{Partie B -- Étude de \(f(x)=\dfrac{\ln x}{x-1}\)}

\begin{enumerate}[label=\textbf{1.\arabic*}, resume, leftmargin=1.6em]

\item \textbf{Signe de \(f(x)\)}

Sur \(]0,1[\), on a
\[
\ln x<0
\qquad \text{et} \qquad
x-1<0,
\]
donc
\[
f(x)>0.
\]

Sur \(]1,+\infty[\), on a
\[
\ln x>0
\qquad \text{et} \qquad
x-1>0,
\]
donc
\[
f(x)>0.
\]

Finalement,
\[
\boxed{f(x)>0 \text{ pour tout } x>0,\ x\neq 1.}
\]

\item \textbf{Calcul de \(f'(x)\)}

\begin{rappel}
Si
\[
f(x)=\frac{u(x)}{v(x)},
\]
alors
\[
f'(x)=\frac{u'(x)v(x)-u(x)v'(x)}{v(x)^2}.
\]
\end{rappel}

Ici
\[
u(x)=\ln x,\qquad u'(x)=\frac1x,
\qquad
v(x)=x-1,\qquad v'(x)=1.
\]

Donc
\[
f'(x)=\frac{\frac1x(x-1)-\ln x}{(x-1)^2}.
\]

Ainsi
\[
\boxed{
f'(x)=\frac{\frac{x-1}{x}-\ln x}{(x-1)^2}.
}
\]

\item \textbf{Étudier le signe de \(\dfrac{x-1}{x}-\ln x\)}

Posons
\[
h(x)=\frac{x-1}{x}-\ln x=1-\frac1x-\ln x.
\]

Alors
\[
h'(x)=\frac1{x^2}-\frac1x=\frac{1-x}{x^2}.
\]

Comme \(x^2>0\), le signe de \(h'(x)\) est celui de \(1-x\). Donc :
\begin{itemize}
    \item \(h'(x)>0\) sur \(]0,1[\),
    \item \(h'(1)=0\),
    \item \(h'(x)<0\) sur \(]1,+\infty[\).
\end{itemize}

Ainsi, \(h\) admet un maximum en \(x=1\). Or
\[
h(1)=1-1-\ln 1=0.
\]

Donc, pour tout \(x>0\),
\[
\boxed{\frac{x-1}{x}-\ln x\le 0,}
\]
avec égalité si et seulement si \(x=1\).

\item \textbf{Variations de \(f\)}

Le dénominateur \((x-1)^2\) est strictement positif pour \(x\neq 1\).  
Le signe de \(f'(x)\) est donc celui de
\[
\frac{x-1}{x}-\ln x,
\]
qui est toujours négatif pour \(x\neq 1\).

Ainsi :
\[
\boxed{f \text{ est strictement décroissante sur } ]0,1[ \text{ et sur } ]1,+\infty[.}
\]

\item \textbf{Calcul de \(\lim\limits_{x\to 1}\dfrac{\ln x}{x-1}\)}

\begin{rappel}
\[
\lim_{x\to 1}\frac{\ln x-\ln 1}{x-1}
\]
est le nombre dérivé de la fonction \(\ln\) en \(1\), soit
\[
(\ln)'(1)=1.
\]
\end{rappel}

Donc
\[
\boxed{\lim_{x\to 1}\frac{\ln x}{x-1}=1.}
\]

\item \textbf{Montrer que \(\ln x\le x-1\)}

D’après la partie A,
\[
x-\ln x\ge 1.
\]

En réarrangeant :
\[
-\ln x \ge 1-x
\qquad\Longleftrightarrow\qquad
\ln x \le x-1.
\]

Donc
\[
\boxed{\ln x\le x-1 \text{ pour tout } x>0.}
\]

\item \textbf{Cas d’égalité}

L’égalité
\[
\ln x=x-1
\]
équivaut à
\[
x-\ln x=1.
\]

Or cela n’a lieu que pour \(x=1\). Donc
\[
\boxed{\text{égalité si et seulement si } x=1.}
\]

\end{enumerate}

\subsection*{Partie C -- Exploitation}

\begin{enumerate}[label=\textbf{1.\arabic*}, resume, leftmargin=1.6em]

\item \textbf{Étudier l’assertion \(\ln x\le \sqrt{x}-1\)}

Posons \(x=t^2\) avec \(t>0\). Alors
\[
\ln x=\ln(t^2)=2\ln t.
\]

L’inégalité devient
\[
2\ln t\le t-1.
\]

Considérons
\[
\varphi(t)=t-1-2\ln t.
\]

On calcule :
\[
\varphi'(t)=1-\frac{2}{t}=\frac{t-2}{t}.
\]

Comme \(t>0\), on en déduit :
\begin{itemize}
    \item \(\varphi'(t)<0\) sur \(]0,2[\),
    \item \(\varphi'(2)=0\),
    \item \(\varphi'(t)>0\) sur \(]2,+\infty[\).
\end{itemize}

Donc \(\varphi\) admet un minimum en \(t=2\). Or
\[
\varphi(2)=2-1-2\ln 2=1-2\ln 2.
\]

Comme
\[
2\ln 2\approx 1{,}386>1,
\]
on a
\[
\varphi(2)<0.
\]

Ainsi l’inégalité n’est pas vraie pour tout \(t>0\), donc pas pour tout \(x>0\).

Un contre-exemple simple est \(x=4\) :
\[
\ln 4\approx 1{,}386
\qquad\text{et}\qquad
\sqrt4-1=1.
\]
Donc
\[
\ln 4>\sqrt4-1.
\]

\[
\boxed{\text{L’assertion est fausse.}}
\]

\begin{reponseattendue}
Un contre-exemple correct et clairement justifié suffisait. Le choix \(x=4\) est le plus naturel.
\end{reponseattendue}

\item \textbf{Étudier le nombre de solutions de \(\ln x=\dfrac{x-1}{2}\)}

Posons
\[
\psi(x)=\ln x-\frac{x-1}{2},
\qquad x>0.
\]

L’équation cherchée est
\[
\psi(x)=0.
\]

Calculons la dérivée :
\[
\psi'(x)=\frac1x-\frac12=\frac{2-x}{2x}.
\]

Comme \(x>0\), on en déduit :
\begin{itemize}
    \item \(\psi'(x)>0\) sur \(]0,2[\),
    \item \(\psi'(2)=0\),
    \item \(\psi'(x)<0\) sur \(]2,+\infty[\).
\end{itemize}

Donc \(\psi\) croît sur \(]0,2]\), puis décroît sur \([2,+\infty[\).

Calculons les limites :
\[
\lim_{x\to 0^+}\psi(x)=-\infty
\qquad\text{car}\qquad
\ln x\to -\infty,
\]
et
\[
\lim_{x\to +\infty}\psi(x)=-\infty
\qquad\text{car}\qquad
\frac{x-1}{2}\to +\infty \text{ domine } \ln x.
\]

Calculons quelques valeurs :
\[
\psi(1)=\ln 1-\frac{1-1}{2}=0.
\]
Donc \(x=1\) est une solution.

De plus,
\[
\psi(2)=\ln 2-\frac12>0.
\]

Ainsi, la fonction part de \(-\infty\), passe par \(0\) en \(x=1\), devient positive, atteint un maximum en \(x=2\), puis redescend vers \(-\infty\). Elle coupe donc l’axe des abscisses une deuxième fois après \(2\).

Conclusion :
\[
\boxed{\text{L’équation admet exactement deux solutions sur } ]0,+\infty[.}
\]

L’une est
\[
\boxed{x=1,}
\]
et l’autre est un réel \(\alpha>2\).

\end{enumerate}

\vspace{0.6em}
\hrule
\vspace{0.6em}

\section*{Exercice 2 -- Probabilités, suite récurrente, vitesse de convergence}

On note \(A_n\) l’événement « recevoir une réduction au \(n\)-ième jour », et
\[
a_n=P(A_n).
\]

On sait que :
\[
a_1=\frac12,
\qquad
P(A_{n+1}\mid A_n)=\frac13,
\qquad
P(A_{n+1}\mid \overline{A_n})=\frac14.
\]

\subsection*{Partie A -- Mise en équation}

\begin{enumerate}[label=\textbf{2.\arabic*}, leftmargin=1.6em]

\item \textbf{Arbre de probabilités}

L’arbre attendu est :
\begin{itemize}
    \item depuis \(A_n\), on va vers \(A_{n+1}\) avec probabilité \(\frac13\), et vers \(\overline{A_{n+1}}\) avec probabilité \(\frac23\) ;
    \item depuis \(\overline{A_n}\), on va vers \(A_{n+1}\) avec probabilité \(\frac14\), et vers \(\overline{A_{n+1}}\) avec probabilité \(\frac34\).
\end{itemize}

\item \textbf{Relation de récurrence}

Par la formule des probabilités totales,
\[
a_{n+1}=P(A_{n+1})
= P(A_n)P(A_{n+1}\mid A_n)+P(\overline{A_n})P(A_{n+1}\mid \overline{A_n}).
\]

Donc
\[
a_{n+1}=a_n\cdot \frac13 +(1-a_n)\cdot \frac14.
\]

En développant,
\[
a_{n+1}=\frac13 a_n+\frac14-\frac14 a_n
=\left(\frac13-\frac14\right)a_n+\frac14
=\frac1{12}a_n+\frac14.
\]

\[
\boxed{a_{n+1}=\frac1{12}a_n+\frac14}
\]

\end{enumerate}

\subsection*{Partie B -- Étude théorique}

\begin{enumerate}[label=\textbf{2.\arabic*}, resume, leftmargin=1.6em]

\item \textbf{Calcul de \(\ell\)}

On résout
\[
\ell=\frac1{12}\ell+\frac14.
\]

Donc
\[
\ell-\frac1{12}\ell=\frac14
\quad\Longrightarrow\quad
\frac{11}{12}\ell=\frac14
\quad\Longrightarrow\quad
\ell=\frac14\cdot \frac{12}{11}=\frac3{11}.
\]

\[
\boxed{\ell=\frac3{11}}
\]

\item \textbf{Montrer que \((u_n)\) est géométrique}

On pose
\[
u_n=a_n-\ell=a_n-\frac3{11}.
\]

Alors
\[
u_{n+1}=a_{n+1}-\frac3{11}.
\]

Or
\[
a_{n+1}=\frac1{12}a_n+\frac14.
\]

Comme
\[
\frac3{11}=\frac1{12}\cdot \frac3{11}+\frac14,
\]
on obtient
\[
u_{n+1}
=\left(\frac1{12}a_n+\frac14\right)-\left(\frac1{12}\cdot \frac3{11}+\frac14\right)
=\frac1{12}\left(a_n-\frac3{11}\right)
=\frac1{12}u_n.
\]

Donc
\[
\boxed{u_{n+1}=\frac1{12}u_n}
\]
et \((u_n)\) est géométrique.

\item \textbf{Raison et premier terme}

La raison vaut
\[
\boxed{q=\frac1{12}}.
\]

Le premier terme est
\[
u_1=a_1-\frac3{11}=\frac12-\frac3{11}=\frac{11-6}{22}=\frac5{22}.
\]

\[
\boxed{u_1=\frac5{22}}
\]

\item \textbf{Expression de \(a_n\)}

Comme \((u_n)\) est géométrique,
\[
u_n=u_1\left(\frac1{12}\right)^{n-1}.
\]

Donc
\[
a_n=\ell+u_n
=\frac3{11}+\frac5{22}\left(\frac1{12}\right)^{n-1}.
\]

\[
\boxed{
a_n=\frac3{11}+\frac5{22}\left(\frac1{12}\right)^{n-1}
}
\]

\item \textbf{Limite de \((a_n)\)}

Comme
\[
\left(\frac1{12}\right)^{n-1}\to 0,
\]
on en déduit
\[
\boxed{\lim_{n\to+\infty}a_n=\frac3{11}.}
\]

\end{enumerate}

\subsection*{Partie C -- Lecture qualitative}

\begin{enumerate}[label=\textbf{2.\arabic*}, resume, leftmargin=1.6em]

\item \textbf{Sens de variation de \((a_n)\)}

On a
\[
a_n=\frac3{11}+\frac5{22}\left(\frac1{12}\right)^{n-1}.
\]

Le deuxième terme est positif et strictement décroissant. Donc
\[
\boxed{(a_n)\text{ est strictement décroissante.}}
\]

\item \textbf{Position de \(a_n\) par rapport à \(\ell\)}

On calcule :
\[
a_n-\ell=\frac5{22}\left(\frac1{12}\right)^{n-1}>0.
\]

Donc
\[
\boxed{a_n>\ell \text{ pour tout } n\ge 1.}
\]

\item \textbf{Interprétation de \(\ell\)}

\begin{reponseattendue}
\(\ell\) représente la probabilité d’obtenir une réduction à long terme, lorsque le système s’est stabilisé.
\end{reponseattendue}

\end{enumerate}

\subsection*{Partie D -- Questions plus fines}

\begin{enumerate}[label=\textbf{2.\arabic*}, resume, leftmargin=1.6em]

\item \textbf{Trouver le plus petit entier \(n\) tel que \(|a_n-\ell|<10^{-4}\)}

On doit résoudre
\[
\frac5{22}\left(\frac1{12}\right)^{n-1}<10^{-4}.
\]

Testons successivement :

Pour \(n=4\),
\[
\frac5{22}\cdot \frac1{12^3}\approx 1{,}31\times 10^{-4},
\]
ce n’est pas suffisant.

Pour \(n=5\),
\[
\frac5{22}\cdot \frac1{12^4}\approx 1{,}09\times 10^{-5},
\]
c’est bon.

Ainsi
\[
\boxed{n=5}
\]
est le plus petit entier convenable.

\item \textbf{Expression de \(S_n=a_1+\cdots +a_n\)}

On a
\[
S_n=\sum_{k=1}^n a_k
=\sum_{k=1}^n\left(\frac3{11}+\frac5{22}\left(\frac1{12}\right)^{k-1}\right).
\]

Donc
\[
S_n=\frac{3n}{11}+\frac5{22}\sum_{k=0}^{n-1}\left(\frac1{12}\right)^k.
\]

\begin{rappel}
Pour une somme géométrique :
\[
1+q+\cdots+q^{n-1}=\frac{1-q^n}{1-q}
\qquad (q\neq 1).
\]
\end{rappel}

Ici,
\[
\sum_{k=0}^{n-1}\left(\frac1{12}\right)^k
=\frac{1-\left(\frac1{12}\right)^n}{1-\frac1{12}}
=\frac{1-\left(\frac1{12}\right)^n}{\frac{11}{12}}
=\frac{12}{11}\left(1-\left(\frac1{12}\right)^n\right).
\]

Ainsi
\[
S_n=\frac{3n}{11}+\frac5{22}\cdot \frac{12}{11}\left(1-\left(\frac1{12}\right)^n\right).
\]

Or
\[
\frac5{22}\cdot \frac{12}{11}=\frac{30}{121}.
\]

Donc
\[
\boxed{
S_n=\frac{3n}{11}+\frac{30}{121}\left(1-\left(\frac1{12}\right)^n\right)
}
\]

\item \textbf{Interprétation de \(S_n\)}

\begin{reponseattendue}
\(S_n\) représente le nombre moyen cumulé de jours avec réduction parmi les \(n\) premiers jours.
\end{reponseattendue}

\end{enumerate}

\vspace{0.6em}
\hrule
\vspace{0.6em}

\section*{Exercice 3 -- Nombres complexes et géométrie}

On considère
\[
z_A=1,\qquad z_B=\ii,\qquad z_C=-1.
\]

\subsection*{Partie A -- Lecture géométrique}

\begin{enumerate}[label=\textbf{3.\arabic*}, leftmargin=1.6em]

\item \textbf{Placement des points}

Les affixes correspondent aux coordonnées :
\[
A(1,0),\qquad B(0,1),\qquad C(-1,0).
\]

\item \textbf{Calcul des longueurs}

\[
AB=|z_B-z_A|=|\ii-1|=|-1+\ii|=\sqrt{(-1)^2+1^2}=\sqrt2.
\]

\[
BC=|z_C-z_B|=|-1-\ii|=\sqrt{1^2+1^2}=\sqrt2.
\]

\[
AC=|z_A-z_C|=|1-(-1)|=2.
\]

Donc
\[
\boxed{AB=\sqrt2,\quad BC=\sqrt2,\quad AC=2.}
\]

\item \textbf{Nature du triangle}

On a
\[
AB^2+BC^2=2+2=4=AC^2.
\]

Donc \(ABC\) est rectangle en \(B\).  
De plus, \(AB=BC\), donc il est isocèle en \(B\).

\[
\boxed{ABC \text{ est rectangle isocèle en } B.}
\]

\end{enumerate}

\subsection*{Partie B -- Transformation \(T\)}

\[
z'=(1-\ii)z+1+\ii.
\]

\begin{enumerate}[label=\textbf{3.\arabic*}, resume, leftmargin=1.6em]

\item \textbf{Module et argument de \(1-\ii\)}

\[
|1-\ii|=\sqrt{1^2+(-1)^2}=\sqrt2.
\]

Un argument de \(1-\ii\) est
\[
-\frac{\pi}{4}.
\]

Donc
\[
\boxed{1-\ii=\sqrt2\,\e^{-\ii\pi/4}.}
\]

\item \textbf{Interprétation géométrique de \(z\mapsto (1-\ii)z\)}

Multiplier par un complexe de module \(\sqrt2\) et d’argument \(-\frac{\pi}{4}\) revient à effectuer :
\begin{itemize}
    \item une homothétie de centre \(O\) et de rapport \(\sqrt2\),
    \item suivie d’une rotation de centre \(O\) et d’angle \(-\frac{\pi}{4}\).
\end{itemize}

\item \textbf{Point fixe}

On cherche \(\omega\) tel que
\[
\omega=(1-\ii)\omega+1+\ii.
\]

Donc
\[
\omega-(1-\ii)\omega=1+\ii.
\]

On factorise :
\[
\ii\,\omega=1+\ii.
\]

Ainsi
\[
\omega=\frac{1+\ii}{\ii}.
\]

Or
\[
\frac1{\ii}=-\ii,
\]
donc
\[
\omega=(1+\ii)(-\ii)=1-\ii.
\]

\[
\boxed{\omega=1-\ii}
\]

\item \textbf{Nature de \(T\)}

Une application de la forme
\[
z'=az+b
\qquad (a\neq 0,\ a\neq 1)
\]
est une similitude directe de centre le point fixe.

Ici
\[
a=1-\ii \neq 0 \quad \text{et} \quad a\neq 1.
\]

Donc
\[
\boxed{T \text{ est une similitude directe de centre } \Omega \text{ d’affixe } 1-\ii.}
\]

\item \textbf{Rapport et angle}

Le rapport vaut
\[
\boxed{\sqrt2}
\]
et l’angle vaut
\[
\boxed{-\frac{\pi}{4}}.
\]

\end{enumerate}

\subsection*{Partie C -- Itérations}

On définit
\[
z_0=2,\qquad z_{n+1}=(1-\ii)z_n+1+\ii.
\]

\begin{enumerate}[label=\textbf{3.\arabic*}, resume, leftmargin=1.6em]

\item \textbf{Calcul de \(z_1\) et \(z_2\)}

\[
z_1=(1-\ii)\cdot 2+1+\ii=2-2\ii+1+\ii=3-\ii.
\]

Donc
\[
\boxed{z_1=3-\ii.}
\]

Puis
\[
z_2=(1-\ii)(3-\ii)+1+\ii.
\]

Développons :
\[
(1-\ii)(3-\ii)=3-\ii-3\ii+\ii^2=3-4\ii-1=2-4\ii.
\]

Donc
\[
z_2=2-4\ii+1+\ii=3-3\ii.
\]

\[
\boxed{z_2=3-3\ii.}
\]

\item \textbf{Relation avec le point fixe}

Comme \(\omega\) vérifie
\[
\omega=(1-\ii)\omega+1+\ii,
\]
on a
\[
z_{n+1}-\omega=(1-\ii)z_n+1+\ii-\big((1-\ii)\omega+1+\ii\big).
\]

Ainsi
\[
\boxed{z_{n+1}-\omega=(1-\ii)(z_n-\omega).}
\]

\item \textbf{Expression de \(z_n-\omega\)}

La relation précédente montre par récurrence que
\[
z_n-\omega=(1-\ii)^n(z_0-\omega).
\]

Or
\[
z_0-\omega=2-(1-\ii)=1+\ii.
\]

Donc
\[
\boxed{z_n-\omega=(1-\ii)^n(1+\ii).}
\]

\item \textbf{Calcul du module}

On prend le module :
\[
|z_n-\omega|=|1-\ii|^n\,|1+\ii|.
\]

Or
\[
|1-\ii|=\sqrt2
\qquad\text{et}\qquad
|1+\ii|=\sqrt2.
\]

Donc
\[
|z_n-\omega|=(\sqrt2)^n|z_0-\omega|.
\]

\[
\boxed{|z_n-\omega|=(\sqrt2)^n|z_0-\omega|.}
\]

\item \textbf{Interprétation géométrique}

\begin{reponseattendue}
La distance à \(\Omega\) est multipliée à chaque étape par \(\sqrt2>1\), donc elle augmente.  
En même temps, l’angle tourne de \(-\frac{\pi}{4}\) à chaque itération.  
Les points \(M_n\) décrivent donc une spirale divergente autour de \(\Omega\).
\end{reponseattendue}

\end{enumerate}

\subsection*{Partie D -- Lieu géométrique}

On considère
\[
\mathcal E=\left\{M(z)\in\C,\ z\neq 1,\ \left|\frac{z-\ii}{z-1}\right|=\sqrt2\right\}.
\]

\begin{enumerate}[label=\textbf{3.\arabic*}, resume, leftmargin=1.6em]

\item \textbf{Interprétation géométrique}

On a
\[
\left|\frac{z-\ii}{z-1}\right|=\frac{|z-\ii|}{|z-1|}.
\]

Donc la condition équivaut à
\[
|z-\ii|=\sqrt2\,|z-1|.
\]

Cela signifie
\[
MB=\sqrt2\,MA.
\]

\[
\boxed{\mathcal E=\{M\mid MB=\sqrt2\,MA\}.}
\]

\item \textbf{Montrer que \(\mathcal E\) est un cercle}

Posons
\[
z=x+\ii y.
\]

Alors
\[
|z-\ii|^2=x^2+(y-1)^2,
\qquad
|z-1|^2=(x-1)^2+y^2.
\]

La condition devient
\[
x^2+(y-1)^2=2\big((x-1)^2+y^2\big).
\]

Développons :
\[
x^2+y^2-2y+1=2x^2-4x+2+2y^2.
\]

En regroupant :
\[
0=x^2+y^2-4x+2y+1.
\]

Complétons les carrés :
\[
x^2-4x=(x-2)^2-4,
\qquad
y^2+2y=(y+1)^2-1.
\]

Donc
\[
(x-2)^2-4+(y+1)^2-1+1=0,
\]
soit
\[
(x-2)^2+(y+1)^2=4.
\]

Ainsi
\[
\boxed{\mathcal E \text{ est le cercle de centre }(2,-1)\text{ et de rayon }2.}
\]

\item \textbf{Déterminer si \(C\in \mathcal E\)}

Le point \(C\) a pour coordonnées \((-1,0)\).

On teste dans l’équation du cercle :
\[
(-1-2)^2+(0+1)^2=9+1=10.
\]

Ce n’est pas égal à \(4\), donc
\[
\boxed{C\notin \mathcal E.}
\]

\end{enumerate}

\vspace{0.6em}
\hrule
\vspace{0.6em}

\section*{Exercice 4 -- Arithmétique, congruences, structure des entiers}

\subsection*{Partie A -- Étude modulo 7}

\begin{enumerate}[label=\textbf{4.\arabic*}, leftmargin=1.6em]

\item \textbf{Restes successifs de \(3^n\) modulo \(7\)}

\[
3^1\equiv 3 \pmod 7,
\]
\[
3^2=9\equiv 2 \pmod 7,
\]
\[
3^3\equiv 3\cdot 2=6 \pmod 7,
\]
\[
3^4\equiv 3\cdot 6=18\equiv 4 \pmod 7,
\]
\[
3^5\equiv 3\cdot 4=12\equiv 5 \pmod 7,
\]
\[
3^6\equiv 3\cdot 5=15\equiv 1 \pmod 7.
\]

\item \textbf{En déduire \(3^6\equiv 1\pmod 7\)}

Le calcul précédent donne directement
\[
\boxed{3^6\equiv 1 \pmod 7.}
\]

\item \textbf{Résoudre \(3^n\equiv 5 \pmod 7\)}

On a vu que
\[
3^5\equiv 5\pmod 7.
\]

Comme la suite des puissances de \(3\) modulo \(7\) est périodique de période \(6\), on obtient
\[
\boxed{3^n\equiv 5\pmod 7 \Longleftrightarrow n\equiv 5\pmod 6.}
\]

\end{enumerate}

\subsection*{Partie B -- Équation \(3x-7y=5\)}

\begin{enumerate}[label=\textbf{4.\arabic*}, resume, leftmargin=1.6em]

\item \textbf{Une solution particulière}

Cherchons une solution simple. Si \(y=1\), alors
\[
3x-7=5
\quad\Longrightarrow\quad
3x=12
\quad\Longrightarrow\quad
x=4.
\]

Donc une solution particulière est
\[
\boxed{(4,1).}
\]

\item \textbf{Toutes les solutions entières}

L’équation homogène associée est
\[
3x-7y=0.
\]

Comme \(3\) et \(7\) sont premiers entre eux, on obtient
\[
x=7k,\qquad y=3k
\qquad (k\in\Z).
\]

Donc toutes les solutions de l’équation complète sont
\[
\boxed{x=4+7k,\qquad y=1+3k,\qquad k\in\Z.}
\]

\item \textbf{Solutions avec \(0\le x\le 20\)}

On impose
\[
0\le 4+7k\le 20.
\]

Donc
\[
-4\le 7k\le 16
\quad\Longrightarrow\quad
-\frac47\le k\le \frac{16}{7}.
\]

Ainsi
\[
k\in\{0,1,2\}.
\]

Les solutions cherchées sont donc :
\[
\boxed{(4,1),\ (11,4),\ (18,7).}
\]

\end{enumerate}

\subsection*{Partie C -- Carrés modulo 8}

\begin{enumerate}[label=\textbf{4.\arabic*}, resume, leftmargin=1.6em]

\item \textbf{Restes possibles de \(n^2\) modulo 8}

Il suffit de tester tous les restes possibles modulo \(8\) :
\[
0^2\equiv 0,\quad 1^2\equiv 1,\quad 2^2\equiv 4,\quad 3^2\equiv 1,
\]
\[
4^2\equiv 0,\quad 5^2\equiv 1,\quad 6^2\equiv 4,\quad 7^2\equiv 1 \pmod 8.
\]

Donc les seuls restes possibles sont
\[
\boxed{0,\ 1,\ 4.}
\]

\item \textbf{Montrer que \(n^2\equiv 3\pmod 8\) est impossible}

Comme \(3\notin \{0,1,4\}\), la congruence
\[
n^2\equiv 3\pmod 8
\]
n’a pas de solution entière.

\[
\boxed{n^2\equiv 3\pmod 8 \text{ n’a aucune solution dans } \Z.}
\]

\item \textbf{Résoudre \(n^2\equiv 1\pmod 8\)}

D’après le tableau précédent, cela se produit exactement lorsque \(n\) est impair.

\[
\boxed{n^2\equiv 1\pmod 8 \Longleftrightarrow n \text{ est impair}.}
\]

\end{enumerate}

\subsection*{Partie D -- Synthèse : \(x^2+7y^2=2026\)}

\begin{enumerate}[label=\textbf{4.\arabic*}, resume, leftmargin=1.6em]

\item \textbf{Pourquoi l’étude modulo \(8\) est pertinente}

Le second membre vaut
\[
2026.
\]

Or
\[
2026=8\times 253+2,
\]
donc
\[
2026\equiv 2\pmod 8.
\]

Comme on connaît précisément les valeurs possibles d’un carré modulo \(8\), travailler modulo \(8\) est naturel et efficace.

\[
\boxed{\text{Le modulo }8\text{ est pertinent car les carrés y ont très peu de valeurs possibles.}}
\]

\item \textbf{Montrer que l’équation n’a pas de solution}

On réduit l’équation
\[
x^2+7y^2=2026
\]
modulo \(8\).

Comme
\[
2026\equiv 2\pmod 8
\qquad\text{et}\qquad
7\equiv -1\pmod 8,
\]
on obtient
\[
x^2-y^2\equiv 2\pmod 8.
\]

Or, d’après la partie C, \(x^2\) et \(y^2\) appartiennent à \(\{0,1,4\}\) modulo \(8\).

Testons toutes les différences possibles :
\[
0-0=0,\quad 0-1\equiv 7,\quad 0-4\equiv 4,
\]
\[
1-0=1,\quad 1-1=0,\quad 1-4\equiv 5,
\]
\[
4-0=4,\quad 4-1=3,\quad 4-4=0 \pmod 8.
\]

Aucune de ces valeurs n’est congrue à \(2\) modulo \(8\).

Donc l’équation est impossible dans \(\Z^2\).

\[
\boxed{x^2+7y^2=2026 \text{ n’admet aucune solution dans } \Z^2.}
\]

\end{enumerate}

\vspace{0.8em}
\begin{center}
\textbf{Fin du corrigé}
\end{center}

\end{document}